\section{Conclusion}
We introduce human and mouse benchmarks that compare peptides from the same parent protein under the same MHC restriction and predict which has stronger presentation evidence in a specified tissue. Peptide sequence supports this ranking in both species, and removing tissue identity substantially reduces performance. Performance also decreases when every held-out peptide is excluded from all fitting combinations, showing that repeated peptide identities make estimates in familiar contexts more optimistic. The retained signal supports tissue-associated peptide ranking in previously represented tissue--MHC combinations but does not identify a specific antigen-processing mechanism. Complex architectures do not consistently outperform strong simpler controls after peptide separation. The main contribution is therefore a biologically focused and reproducible benchmark that distinguishes tissue-associated signal from general peptide--MHC compatibility while making the limits of unseen-peptide prediction explicit. Validation in unseen tissues, alleles, studies, and prospective samples remains necessary before biological or clinical application.
